\newpage
%

\chapter{Compiling, running and debugging}
\label{sec_compilationrunning}

\section{Compiling OASIS3-MCT}
\label{subsec_compile}

OASIS3-MCT is an MPI code. Compiling OASIS3-MCT libraries can be done
from the {\tt oasis3-MCT/util/make\_dir} directory
with the makefile {\tt TopMakefileOasis3}.
{\tt TopMakefileOasis3} includes the header file {\tt make.inc}
which should then point to (include) your own {\tt make.your\_platform}
file.  That file is specific to the hardware and compiling platform used.

Several header files are distributed with the release and can by used as
a template to create a custom file for your machine.
The root of the OASIS3-MCT tree can be anywhere, but it must be defined
by the variable {\tt COUPLE}.  Similarly, the variable {\tt ARCHDIR} defines
the location of the compilation directory.
Both of these variables are usually defined in the
{\tt make.your\_platform} file.  Finally, the OASIS3-MCT library
should be compiled with the same compilers and system software as any
coupled models using it.  After successful compilation, resulting
libraries are found in the directory in {\tt \$ARCHDIR/lib} while
object and module files are found in {\tt \$ARCHDIR/build} and {\tt \$ARCHDIR/build\_shared}. 

OASIS3-MCT has historically created static libraries for use in
Fortran source codes.  However, C language bindings are now
available, and python codes are now fully supported. Therefore, the OASIS3-MCT makefile {\tt TopMakefileOasis3}
supports compilation of both static and shared libraries.

{\tt TopMakefileOasis3} has several targets such as:

\begin{itemize}
\item  {\tt oasis3-psmile   }  = static oasis libraries (this is the default)
\item  {\tt static-libs     }  = static oasis libraries (same as oasis3-psmile)
\item  {\tt static-cbindings}  = static oasis libraries plus static c-bindings
\item  {\tt shared-libs     }  = shared (dynamic) oasis libraries
\item  {\tt shared-cbindings}  = shared oasis libraries plus shared c-bindings
\item  {\tt pyoasis         }  = builds and installs shared-cbindings package and installs files containing higher and intermediate python classes for use with python codes
\item  {\tt realclean       }  = cleans and resets the build
\end{itemize}

The names of the libraries produced
are {\it mct}, {\it mpeu},  {\it scrip}, {\it psmile.MPI1}, and {\it oasis.cbind}
with standard prefixes ({\it lib}) and suffixes ({\it .a} or {\it .so}).

The following targets have been used historically to compile oasis for Fortran codes:

\begin{itemize}
\item {\tt make -f TopMakefileOasis3 help}

  provides a current list of available targets.

\item {\tt make -f TopMakefileOasis3 realclean}

  removes all OASIS3-MCT compiled sources and librairies.

\item {\tt make -f TopMakefileOasis3} or

      {\tt make -f TopMakefileOasis3 oasis3\_psmile}

  compiles static versions of OASIS3-MCT Fortran libraries {\it mct}, {\it mpeu},
  {\it scrip} and {\it psmile};

\end{itemize}

Log and error messages from compilation are normally saved in the directory
{\tt /util/make\_dir} in the files
{\tt COMP.log} and {\tt COMP.err} or similar.  The {\tt TopMakefileOasis3}
output will direct users to the compile output files.

To interface a component code with OASIS3-MCT and use {\tt mod\_oasis} (see section \ref{subsubsec_Use}), it is required to include OASIS3-MCT modules from {\tt \$ARCHDIR/include} and link with appropriate libraries in {\tt \$ARCHDIR/lib} during the compilation and linking.

Exchange of coupling fields in single and double precision is now supported directly through the interface 
(see section \ref{subsubsec_Declaration}).  Single precision fields are converted to double precision fields internally and temporarily.
For double precision coupling fields, there is no need to promote REAL variables to DOUBLE PRECISION at compilation; this is done automatically within the OASIS3-MCT library.

\section{CPP keys}
\label{subsec_cpp}

The following OASIS3-MCT CPP keys can be specified in {\tt CPPDEF} in {\tt make.{\it your\_platform}} file:
\begin{itemize}
\item {\tt TREAT\_OVERLAY}:  ensures, in {\tt SCRIPR/CONSERV}
  remapping (see section \ref{subsec_interp}), that if two cells of
  the source grid overlay and none is masked a priori, the one with the greater numerical
  index will not be considered (they also can be both masked); this is mandatory
  for this remapping. For example, if the grid line with {\it i=1} overlaps
  the grid line with {\it i=imax}, it is the latter that must be masked;
  when this is not the case with the mask defined in {\it masks.nc},
  this CPP key forces these rules to be respected.

\item {\tt \_\_NO\_16BYTE\_REALS}:  {\bf must} be specified  if you compile
with {\bf PGF90}.
\end{itemize}

%\item {\tt balance}: Add a MPI\_Wtime() function before and after
%  mct\_isend (MPI put) and mct\_recv (MPI get) to calculate the time
%  of the send and receive of a coupling field. This option can be used
%  to produce timestamps in OASIS3-MCT debug files. During a post-processing
%  phase, this information can be used to perform an analysis of the
%  coupled components load (un)balance; specific tools have been
%  developed to do this and will be released with a further version of
%  OASIS3-MCT. {\bf This option is temporarily not recommended as it was observed that
%  it was increasing the simulation time of coupled models on
%  the PRACE computer MareNostrum.}

\section{Running OASIS3-MCT}
\label{subsec_running}

Examples of running environments are provided with the sources. See also the instructions on OASIS web site at
{\tt
  https://verc.enes.org/oasis/oasis-dedicated-user-support-1/user-toys}
for more details.

\subsection{The tutorial toy model}
\label{subsec_tutorial}

A practical example  is provided in {\tt
  oasis3-mct/examples/tutorial} reproducing the coupling between two
simple codes, model1 and model2, with OASIS3-MCT. 

See
the document {\tt tutorial.pdf} there in and the on-line
tutorial at \\
{\tt
  https://verc.enes.org/oasis/oasis-dedicated-user-support-1/user-toys}

{\tt
 /tutorial-and-test\_interpolation-of-oasis3-mct-1}
for more details.

\subsection{The test\_1bin\_ocnice toy model}
\label{subsec_1bin_ocnice}

Another practical example on how to use OASIS3-MCT is provided in \\ {\tt oasis3-mct/examples/test\_1bin\_ocnice}. 
This toy model reproduces the coupling between 5 sub-components within
the “ocnice” executable. See
the document {\tt test\_1bin\_ocnice.pdf} there in for more details.

\subsection{The test\_interpolation environment}
\label{subsec_testinterpolation}

This environment available with the sources in {\tt
  oasis3-mct/examples/test\_interpolation} allows the user to test the
quality of the interpolations and transformations between his source
and target grids by calculating the error of interpolation on the
target grid. For more details, see also the document {\tt test\_interpolation.pdf} there in.

\section{Debugging}
\label{subsec_debug}

\subsection{Debug files}
If you experience problems while running your coupled model with
OASIS3-MCT, you can obtain more information on what is happening by
increasing the {\tt \$NLOGPRT} value in your {\it namcouple}, see section
\ref{subsec_namcouplefirst} for details.

\subsection{Time statistics files}
\label{timestat}

The variable TIMER\_Debug, defined in the {\it namcouple} (second
number on the line below \$NLOGPRT keyword), is used to obtain time
statistics over all the processors for each routine.

Different output are written (in files named {\tt *.timers\_xxxx})
depending on TIMER\_Debug value :

\begin{itemize}
\item {TIMER\_Debug=0} : nothing is calculated, nothing is written.
\item {TIMER\_Debug=1} : the times are calculated and written in a
  single file by the process 0 as well as the min and the max times
  over all the processes.
\item {TIMER\_Debug=2} : the times are calculated and each process
  writes its own file ; process 0 also writes the min and the max
  times over all the processes in its file.
\item {TIMER\_Debug=3} : the times are calculated and each process
  writes its own file ; process 0 also writes in its file the min
  and the max times over all processes and also writes in its file
  all the results for each process.
\end{itemize}

Note that TIMER\_Debug can also be set to -1 to activate the {\it lucia}
tool that can be used to perform an analysis of the coupled components
load balance. More information can be found in the README file in {\tt
  oasis3-mct/util/lucia} directory and report mentioned therein.
 
The time given for each timer is calculated by the difference between
calls to {\tt oasis\_timer\_start()} and {\tt oasis\_timer\_stop()}
and is the accumulated time over the entire run. Here is an overview
of the meaning of the different timers as implemented by default.
\footnote{Many other measures can be obtained by defining the logical
{\tt local\_timers\_on} as {\tt .true.} in different routines or by
implementing other timers. Of course, OASIS3\_MCT and the model code then have to be recompiled.}

\begin{itemize}

\item {'total'} : total time of the simulation, implemented
  in {\tt mod\_oasis\_method}, i.e. between the end of {\tt
    oasis\_init\_comp} and the {\tt
    mpi\_finalize} in routine {\tt oasis\_terminate}.

\item {'init\_thru\_enddef'} : time between the end of {\tt
    oasis\_init\_comp} and the end of {\tt oasis\_enddef}, implemented
  in {\tt mod\_oasis\_method}.

\item {'part\_definition'} : time spent in routine {\tt oasis\_def\_partition}.

\item {'oasis\_enddef'} : time spent in 
  routine {\tt oasis\_enddef}; this routine performs basically all the
  important steps to initialize the coupling exchanges, e.g. the
  internal management of the partition and variable definition, the
  definition of the patterns of communication between the source and
  target processes, the reading of the remapping weight-and-address
  file and the initialisation of the sparse matrix vector multiplication.
\item {'grcv\_00x'} : time spent in the reception of field x in {\tt
    mct\_recv} (including communication and possible waiting time
  linked to unbalance of components).
\item {'wout\_00x'} : time spent in the I/O for field x in routine
  {\tt oasis\_advance\_run}.
\item {'gcpy\_00x'} : time spent in routine {\tt oasis\_advance\_run}
  in copying the field x just received in a local array.
\item {'pcpy\_00x'} : time spent in routine {\tt oasis\_advance\_run}
  in copying the local field x in the array to send (i.e. with local
  transformation besides division for averaging).
\item {'pavg\_00x'} : time spent in routine {\tt oasis\_advance\_run}
  to calculate the average of field x (if done).
\item {'pmap\_00x'/'gmap\_00x'} : time spent in routine {\tt
    oasis\_advance\_run} for the matrix vector multiplication for
  field x on the source/target processes.
\item {'psnd\_00x'} : time spent in routine {\tt oasis\_advance\_run}
  for sending field x (i.e. including call to {\tt mct\_waitsend} and
  {\tt mct\_isend}).
\item {'wtrn\_00x'} : time spent in routine {\tt oasis\_advance\_run}
  to write fields associated with non-instant loctrans operations to
  restart files  (see section \ref{subsec_restartdata} for details).
\item {'wrst\_00x'} : time spent in routine {\tt oasis\_advance\_run}
  to write fields to
  restart files (see section \ref{subsec_restartdata} for details).
\end{itemize}

